\documentclass[12pt,a4paper]{article}
\usepackage[T1]{fontenc}
\usepackage[utf8]{inputenc}
\usepackage{amsmath, amssymb, amsthm}
\usepackage{geometry}
\geometry{margin=1in}
\usepackage{xcolor}
\usepackage{booktabs}
\usepackage[hidelinks]{hyperref}
\setlength{\parindent}{0pt}
\setlength{\parskip}{1em}

% --- legend colour bullet helper (replaces emoji that break pdflatex) ---
\newcommand{\cbull}[1]{\textcolor{#1}{\rule[0.1ex]{0.9em}{0.9em}}}

\title{Ethical Framework for Invisible Groups \\[4pt]
\large A Corrected Information-Theoretic Model (v3.1)}
\author{Based on the Adoptee Test and Identity Integrity}
\date{\today}

\begin{document}

\maketitle

\begin{abstract}
Many ethical scoring systems assign near-zero weight to invisible groups,
allowing severe harm to go undetected. This report presents a corrected,
internally consistent framework that unifies harm, visibility, identity
continuity, and information loss. The Adoptee Test is formulated as an
omission alarm equal to harm multiplied by identity entropy. Crucially, the
link between detection probability and identity entropy is established by
\emph{definition} rather than assertion: detection probability is defined as
$w(g)=2^{-H_{\mathrm{id}}(g)}$, which makes the alarm
$A(g)=\mathcal{H}(g)\,H_{\mathrm{id}}(g)$ a derived identity, not a claimed
coincidence. Boundary conditions are treated as discrete jumps rather than
derivatives. Identity entropy measures truth lost about origins, and the
Ethical Validity Ratio is shown to be recall (sensitivity). A worked example
and a real-world mapping (sealed birth records) demonstrate applicability.
All information quantities are expressed in bits ($\log_2$) throughout.
\end{abstract}

\section{Notation Reconciliation}

The original sheets reused single symbols for incompatible quantities. The
table below fixes every collision; all later sections use only the v3.1
column.

\begin{center}
\begin{tabular}{lll}
\toprule
\textbf{Old symbol} & \textbf{Problem} & \textbf{v3.1 symbol} \\
\midrule
$H(g)$ (harm) & Collides with entropy $H$ & $\mathcal{H}(g)$ \\
$H$ (entropy) & Collides with harm & $H_{\mathrm{id}}$ (identity entropy, bits) \\
$W(g)$ & Overloaded: weight / visibility / detection & $w(g)\in(0,1]$ (detection prob.) \\
$T,L$ & Declared but never used & Removed; truth lost $=\Delta H_{\mathrm{id}}$ \\
$P(g)$ (power) & Appears only in legend, no equation & Removed \\
\bottomrule
\end{tabular}
\end{center}

\section{Weighted Ethical Score}
\[
E(S) = \sum_{g \in G} w(g)\bigl( R(g) - \mathcal{H}(g) \bigr)
\]
where $R(g)$ is rights retained by group $g$. If $w(g)\to 0$, the group's
suffering does not affect $E(S)$. This is the core defect the remaining
sections detect.

\section{Adoptee Test --- Omission Alarm (Surprisal Form)}
\[
A(g) = \mathcal{H}(g)\,\bigl(-\log_2 w(g)\bigr)
\]
Here $-\log_2 w(g)$ is the information content (surprisal, in bits) of
detecting the group. When $w(g)=1$, the group is fully visible and $A(g)=0$:
full visibility cannot constitute an omission defect. As $w(g)\to 0^{+}$,
$-\log_2 w(g)\to\infty$, so $A(g)\to\infty$.

\noindent\textbf{Interpretation:} the more harm a group experiences while
remaining unseen by the framework, the greater the moral defect in the
\emph{framework itself}, not in the group.

\section{Boundary-Condition Test (Discrete, Multiplicative Identity)}
Define Identity Integrity
\[
I = F \cdot K \cdot C \cdot N, \qquad F,K,C,N \in [0,1],
\]
with $F$ = factual knowledge of origins, $K$ = kinship continuity,
$C$ = cultural identity, $N$ = legal identity continuity. The product form is
deliberate: if \emph{any} factor $\to 0$, then $I \to 0$. An additive form
would let three intact factors mask the severance of a fourth, which is
precisely the evasion this test exists to expose. Baseline (all intact) is
$I_{\text{before}} = 1$.

Adoption is a discrete event, not a smooth decline, so $dI/dt$ is undefined at
the jump. Use the jump itself:
\[
\Delta I = I_{\text{after}} - I_{\text{before}} \le 0 .
\]
A framework passes only if it detects $|\Delta I| > \tau$ for a stated
threshold $\tau$; otherwise Framework Failure $= 1$, with failure magnitude
graded by the undetected $|\Delta I|$.

\section{Information-Theoretic Formulation}
The ground truth about a person's origin is a single fact: a point mass, so
its entropy is $H_{\mathrm{id}}^{*}=0$. The adoptee's induced distribution over
possible origins has
\[
H_{\mathrm{id}} = -\sum_{x} p(x)\,\log_2 p(x),
\]
where $p(x)$ is the probability of each possible true fact $x$. Truth lost is
the entropy increase: $\Delta H_{\mathrm{id}} = H_{\mathrm{id}} - 0 =
H_{\mathrm{id}}$.

\paragraph{The unification (by definition, not assertion).}
A previous draft asserted $H_{\mathrm{id}} = -\log_2 w(g)$ as if it followed
automatically. It does not: a full-distribution entropy and a single-event
surprisal are distinct quantities. We therefore \emph{define} the detection
probability as
\[
\boxed{\,w(g) \equiv 2^{-H_{\mathrm{id}}(g)}\,}
\qquad\Longrightarrow\qquad
-\log_2 w(g) = H_{\mathrm{id}}(g).
\]
This is the probability of recovering the truth in a single draw from a
maximum-entropy origin distribution, which is exactly the ``does the framework
retain the truth'' reading of $w$. Substituting into Section~3 gives a clean
derived identity:
\[
A(g) = \mathcal{H}(g)\,H_{\mathrm{id}}(g)
     = \text{harm} \times \text{identity entropy}.
\]
Sections~3, 5, and~6 are now the same object viewed three ways, by
construction.

\paragraph{Contradiction test.} If a system tolerates large increases in
uncertainty about selfhood ($\Delta H_{\mathrm{id}} > 0$) while claiming to
protect human welfare, it asserts a contradiction.

\section{Ethical Validity Ratio as Recall}
Let TP $=$ rights protected and FN $=$ rights ignored. Then the Ethical
Validity Ratio is exactly recall (sensitivity):
\[
\text{Recall} = \frac{\text{TP}}{\text{TP}+\text{FN}}
= \frac{\text{Rights Protected}}{\text{Rights Protected}+\text{Rights Ignored}}.
\]
Naming it recall locates the construct inside standard detection theory and
invites its companions: precision $=\text{TP}/(\text{TP}+\text{FP})$ and
$F_1 = 2\,\text{TP}/(2\,\text{TP}+\text{FP}+\text{FN})$. When hidden rights
(e.g.\ knowledge of origins) are added back into FN, recall collapses --- which
is the whole point.

\section{Worked Example: Adoptee Group (in bits)}
Let a group of adoptees have standardized harm $\mathcal{H}(g)=100$ and
detection probability $w(g)=0.1$ (10\% visibility). Then identity entropy is
\[
H_{\mathrm{id}}(g) = -\log_2(0.1) = \log_2(10) \approx 3.3219 \text{ bits},
\]
and the omission alarm is
\[
A(g) = \mathcal{H}(g)\,H_{\mathrm{id}}(g)
     = 100 \times 3.3219 \approx 332.19 \text{ (harm-bits)}.
\]
A large alarm, flagging a severe framework defect. Contrast the weighted
score: this group contributes $w(g)\bigl(R(g)-\mathcal{H}(g)\bigr)
= 0.1\,(R(g)-100)$. Even with $R(g)=0$ that is only $-10$ to $E(S)$ ---
trivially outweighed by visible groups. The same invisibility that mutes the
group in $E(S)$ is what drives $A(g)$ toward infinity.

\section{Real-World Mapping: Sealed Birth Records}
Consider a US state that seals original birth records for adopted adults. A
review board might score:
\begin{itemize}
\item Rights protected (TP): partial access to medical history.
\item Rights ignored (FN): knowledge of identity, heritage, biological family.
\end{itemize}
If the board includes no adoptee representatives, modelled detection
probability $w(g)$ is low, so $H_{\mathrm{id}}(g)$ and the omission alarm are
high. When the ignored rights are entered into FN, recall collapses. The
framework therefore predicts that the board's self-assessment will report a
falsely high ethical validity --- an excludable-population artefact, not a
real result. (The board$\to w$ link is a modelling assumption, stated as such,
not an identity.)

\section{Conclusion}
By fixing the notation collisions, replacing the undefined derivative with a
discrete jump, grounding the boundary test in a multiplicative identity, and
\emph{defining} $w(g)=2^{-H_{\mathrm{id}}(g)}$ so the unification is derived
rather than asserted, the framework becomes internally consistent and
resistant to technical dismissal. The Adoptee Test functions as a canary: if a
system cannot detect harm to invisible groups, the system itself is the defect.

\section*{Legend (Colours)}
\begin{itemize}
\item[\cbull{green!60!black}] Green: protection / continuity / retained rights
\item[\cbull{red!80!black}] Red: harm / loss / ignored rights
\item[\cbull{blue!70!black}] Blue: visibility / detection / evaluation ($w(g)$)
\item[\cbull{orange!85!black}] Orange: knowledge / origins ($F$)
\item[\cbull{violet}] Purple: legal / identity continuity ($N$)
\end{itemize}

\end{document}